\documentclass[11pt]{article}

\usepackage[T1]{fontenc}
\usepackage[utf8]{inputenc}
\usepackage{amsmath,amssymb,amsthm}
\usepackage{booktabs}
\usepackage{float}
\usepackage{geometry}
\usepackage{hyperref}
\usepackage{url}

\geometry{margin=1in}

\newtheorem{theorem}{Theorem}
\newtheorem{computationalresult}{Computational Result}
\newtheorem{proposition}{Proposition}
\newtheorem{lemma}{Lemma}
\newtheorem{corollary}{Corollary}
\theoremstyle{definition}
\newtheorem{definition}{Definition}
\theoremstyle{remark}
\newtheorem{remark}{Remark}

\title{Computational Extension of OEIS Sequence A096242:\\
Four New Terms for Base-9 Deletable Primes}
\author{Carlo Corti\\
Independent researcher\\
\texttt{carlo.corti@outlook.com}}
\date{June 2026}

\begin{document}

\maketitle

\begin{abstract}
OEIS sequence A096242 counts \(n\)-digit base-9 deletable primes.  The
local OEIS snapshot used for this project listed the first ten values,
\[
  4,14,58,221,911,3638,14687,61435,262189,1140171.
\]
We report a deterministic bottom-up computation extending the table through
\(n=14\):
\[
  a(11)=5056361,\quad
  a(12)=22726556,\quad
  a(13)=103351642,\quad
  a(14)=475048816.
\]
The computation uses the deletion relation in reverse: every element of
\(D_n\) must be obtainable by inserting one base-9 digit into an element of
\(D_{n-1}\).  This gives the explicit raw insertion-attempt count
\(9n\,a(n-1)\), an upper bound on the post-leading-zero candidate count
\(C_n\), before further filtering, global deduplication, and primality
testing.
The method is not heuristic.  It is supported by replay of the known OEIS
values, segmented sort/merge generation for large levels, SHA-256 manifests
for the saved levels, and a separate saved-level validity check through
\(n=14\).  The latter verifies the saved values generated by the production
run, but is not a second exhaustive enumeration of missing candidates.  As
an additional check, the first new level \(D_{11}\) was reproduced both by
cross-mode production reruns and by a separate GMP-based re-enumerator,
yielding a byte-identical level file.  No corresponding independent
re-enumeration is reported here for \(D_{12}\), \(D_{13}\), or \(D_{14}\).
\end{abstract}

\noindent\textbf{2020 Mathematics Subject Classification.}
Primary 11Y55; Secondary 11A41, 11B83.

\noindent\textbf{Key words and phrases.}
integer sequence, OEIS, deletable prime, base 9, primality testing,
computational number theory, segmented enumeration.

\section{Introduction}
\label{sec:intro}

A deletable prime is a prime whose digit string can be reduced by deleting
digits, one at a time, while remaining within the same recursively defined
class; see the MathWorld overview for the general notion
\cite{WeissteinDeletable}.  OEIS sequence A096242 is the base-9 counting
sequence for this property:
\[
  a(n)=\#\{\text{\(n\)-digit base-9 deletable primes}\}.
\]

The local OEIS snapshot used as source material gives the title ``Number of
\(n\)-digit base-9 deletable primes'', offset \(1,1\), and the first ten
terms \cite{OEIS-A096242}.  It records Michael Kleber as the original
author, dated February 28, 2003; five further terms from Ryan Propper, dated
July 19, 2005; and a Mathematica program by Robert Price, dated November 13,
2018.  The local snapshot was used as the source record for the first ten
terms.  No separate retrieval date for the local snapshot is recorded in the
project artifacts consulted for this manuscript.  The human author checked
the live OEIS entry during manuscript preparation on June 6, 2026, and found
no change relevant to the first ten source terms; no automated live OEIS
check was performed by the AI system for this draft.

The OEIS entry places A096242 in the broader deletable-prime family, with
cross-references including A080603, A080608, and the base-specific sequences
A096235--A096246.  This note focuses only on the base-9 member of that
family.

The OEIS Mathematica program follows a top-down perspective: for each
\(n\), it scans the interval \([9^{n-1},9^n)\), selects primes, and checks
whether deleting one digit reaches a previously stored deletable prime.  The
computation reported here uses the same definition but reverses the search.
Instead of scanning the full interval, it generates candidates only from
\(D_{n-1}\).  The resulting candidate count after leading-zero exclusion and
before the base-9 last-digit filter and deduplication, denoted \(C_n\)
below, satisfies
\[
  C_n \leq 9n\,a(n-1).
\]

The main result is the following finite computational extension.

\begin{computationalresult}[Validated production computation]
\label{thm:main}
Subject to the documented execution of the production C17 implementation and
the recorded validation artifacts described below, the next four values of
A096242 after the local OEIS snapshot are
\[
  a(11)=5056361,\quad
  a(12)=22726556,\quad
  a(13)=103351642,\quad
  a(14)=475048816.
\]
\end{computationalresult}

\begin{corollary}
\label{cor:bfile}
Combining Computational Result~\ref{thm:main} with the ten terms in the
local OEIS source snapshot gives the proposed b-file extension through
\(n=14\) shown in Table~\ref{tab:terms}.
\end{corollary}

The rest of the paper explains the definitions, the exact bottom-up
enumeration, the base-9 filtering, the segmented implementation, and the
validation evidence.  Section~\ref{sec:def} fixes notation.
Section~\ref{sec:bottom-up} proves the bottom-up generation lemma.
Section~\ref{sec:filters} records the base-9 last-digit filter and the role
of global deduplication.  Section~\ref{sec:specificity} records
sequence-specific features of the base-9 search.  Sections~\ref{sec:impl}
and~\ref{sec:validation} describe implementation and validation.
Section~\ref{sec:results} gives the resulting table and the proof of
Computational Result~\ref{thm:main}.

\section{Base-9 deletable primes}
\label{sec:def}

Let \(b=9\), and let \([w]_9\) denote the integer represented by a base-9
digit string \(w\).  A deletion is valid only when the resulting string has
no leading zero.  The recursion is started explicitly from
\(D_1=\{2,3,5,7\}\), which is the finite-set version of the OEIS convention
that a one-digit prime may terminate the deletion process.

Equivalently, an insertion into an \((n-1)\)-digit string is valid for level
\(n\) only if the resulting first digit is nonzero; a deletion from an
\(n\)-digit string is valid only if the resulting \((n-1)\)-digit string is
nonempty and has nonzero first digit.

\begin{definition}
Let
\[
  D_1=\{2,3,5,7\}.
\]
For \(n\geq 2\), an \(n\)-digit base-9 integer \(p=[w]_9\) belongs to
\(D_n\) if \(p\) is prime and there is a valid deletion of one base-9 digit
from \(w\) whose resulting \((n-1)\)-digit integer belongs to \(D_{n-1}\).
Define
\[
  a(n)=|D_n|.
\]
\end{definition}

The deletion relation gives a directed acyclic graph by digit length,
\[
  D_n \longrightarrow D_{n-1}\longrightarrow \cdots \longrightarrow D_1.
\]
This graph is not a tree.  A value may have more than one valid deletion
parent, and the same candidate may be generated by several insertions.
Therefore the enumeration must count distinct values, not insertion paths.
Here ``parent'' means an element of \(D_{n-1}\) whose digit string can be
reached from the value in \(D_n\) by one valid deletion, or equivalently
from which the value in \(D_n\) can be generated by one valid insertion.

\section{Bottom-up generation}
\label{sec:bottom-up}

\begin{lemma}[Bottom-up completeness]
\label{lem:bottom-up}
For \(n\geq 2\), \(D_n\) is exactly the set of primes among the distinct
\(n\)-digit base-9 integers obtained by inserting one digit
\(d\in\{0,\ldots,8\}\) into the base-9 representation of an element of
\(D_{n-1}\), excluding insertions that create a leading zero.
\end{lemma}

\begin{proof}
Let \(p\in D_n\).  By definition, at least one valid deletion of a digit
from \(p\) gives an element \(q\in D_{n-1}\).  Re-inserting the deleted digit
in the deleted position reconstructs \(p\), and the insertion cannot create a
leading zero because \(p\) is an \(n\)-digit integer.  Thus every element of
\(D_n\) appears in the bottom-up candidate set.

Conversely, suppose an \(n\)-digit candidate \(p\) is obtained by inserting
one digit into some \(q\in D_{n-1}\), and suppose \(p\) is prime.  Deleting
the inserted digit gives \(q\), with no leading zero in the resulting string
because \(q\) is a valid \((n-1)\)-digit base-9 integer.  Therefore
\(p\in D_n\).  Passing to distinct generated values removes
multiple-generation ambiguity without changing the set.
\end{proof}

The lemma is the mathematical reason the search is exact rather than a
probabilistic shortcut.  Let \(C_n\) denote the number of candidates
remaining after the leading-zero exclusion and before the base-9 last-digit
filter, deduplication, and primality testing.  This \(C_n\) is a sizing
upper-bound quantity, not the implementation's post-filter or
post-deduplication survivor count.  The raw insertion-attempt count before
that exclusion is exactly \(9n\,a(n-1)\), so
\[
  C_n\le 9n\,a(n-1),
\]
where the factor \(n\) counts the possible insertion positions in the
resulting \(n\)-digit string and the factor \(9\) counts the inserted digit.

\section{Candidate filters and deduplication}
\label{sec:filters}

Base 9 gives a particularly transparent necessary filter.  Since
\(9\equiv 0\pmod 3\), every base-9 integer is congruent modulo \(3\) to its
last base-9 digit.  Therefore a multi-digit prime cannot have final base-9
digit \(0\), \(3\), or \(6\).

\begin{proposition}[Last-digit rejection]
\label{prop:last-digit}
For \(n\geq 2\), any \(p\in D_n\) has final base-9 digit outside
\(\{0,3,6\}\).
\end{proposition}

\begin{proof}
If the final base-9 digit is \(0\), \(3\), or \(6\), then \(p\equiv 0
\pmod 3\).  Such a multi-digit number is greater than \(3\), so it is
composite.
\end{proof}

Thus the only digit-local necessary condition used here is
\[
  \text{last base-9 digit}\in\{1,2,4,5,7,8\}.
\]
There is no additional last-digit parity filter in base \(9\).  Since
\(9\equiv 1\pmod 2\), the parity of a base-9 integer is determined by the
sum of its digits modulo \(2\), not by the last digit alone.

The implementation records the divisibility-by-\(3\) rejection explicitly as
the last-digit filter.  This is a necessary rejection test, not a primality
test.  After global deduplication, the implementation applies small-prime
trial division by the primes
\[
  2,3,5,7,11,13,17,19,23,29,31,37,41,43,47,53,59,61,67,71,73,79,83,89,97,
\]
and then a deterministic Miller--Rabin strong probable-prime test for the
\texttt{uint64\_t} range.  Miller--Rabin is used here in its strong
probable-prime form \cite{Miller1976,Rabin1980}.  The implementation uses
the seven witnesses
\[
  2,\ 325,\ 9375,\ 28178,\ 450775,\ 9780504,\ 1795265022,
\]
which appear in Sinclair's deterministic-base table and are used here as the
implemented seven-witness set for the 64-bit range
\cite{SinclairSPRP}.  The implementation only permits levels up to
\(n=20\), so every tested value is below \(9^{20}<2^{64}\).  The
deterministic sufficiency of this exact seven-witness set for the 64-bit
range rests on Sinclair's computational table, not on a published proof of
this particular composite-base set.  Peer-reviewed background on strong
pseudoprimes and deterministic Miller--Rabin base sets is given by Jaeschke
and by Sorenson--Webster \cite{Jaeschke1993,SorensonWebster2017}; in
particular, Sorenson--Webster gives a peer-reviewed deterministic prime-base
alternative covering the \(2^{64}\) range.  The work of
Fori\v{s}ek--Jan\v{c}ina gives additional machine-word primality-testing
context and discusses faster alternatives to repeated Miller--Rabin rounds
\cite{ForisekJancina2015}.

Global deduplication is indispensable.  Since the deletion graph is a DAG,
not a tree, the same candidate may be generated from different parents or
from different insertion positions.  The value \(a(n)\) is the cardinality of
the set \(D_n\), so repeated insertion paths must not be counted as repeated
deletable primes.

\section{Sequence-specific features of the base-9 search}
\label{sec:specificity}

The contribution here is a finite exact extension, not a proposed density
law for base-9 deletable primes.  There is no closed form for \(a(n)\) in
the project evidence, and no probabilistic model is used to justify the four
new terms.  The useful specificity is instead structural.

First, the recursive definition gives an exact frontier: to build \(D_n\),
one only needs insertions into \(D_{n-1}\).  This is a different search
space from the top-down interval scan in the OEIS program.  Table~\ref{tab:bounds}
shows the raw insertion bounds for the new levels and for the next
uncomputed level.

\begin{table}[ht]
\centering
\begin{tabular}{r r r}
\toprule
Level \(n\) & Parent count \(a(n-1)\) & Bound \(9n\,a(n-1)\)\\
\midrule
11 & 1,140,171 & 112,876,929\\
12 & 5,056,361 & 546,086,988\\
13 & 22,726,556 & 2,659,007,052\\
14 & 103,351,642 & 13,022,306,892\\
15 & 475,048,816 & 64,131,590,160\\
\bottomrule
\end{tabular}
\caption{Raw bottom-up insertion bounds for the computed extension and the
next uncomputed level.}
\label{tab:bounds}
\end{table}

Second, base 9 makes the divisibility-by-\(3\) filter a digit-local
condition: the last base-9 digit alone determines the residue modulo \(3\).
This is a small arithmetic feature, but it makes the sequence less generic
than a base-independent recursive prime enumeration.

Third, the graph has confluent paths.  A candidate that has several valid
deletions, or several insertion histories, is still one element of \(D_n\).
This is why the reproducible computation is naturally phrased in terms of
sort/unique and, at large levels, sorted runs followed by k-way merge.

Finally, the stop at \(n=14\) is a perimeter choice.  From the measured value
\(a(14)=475048816\), the \(n=15\) insertion bound is
\[
  C_{15}\leq 9\cdot 15\cdot 475048816=64131590160.
\]
Stored as \texttt{uint64\_t} values, this is about \(513.1\) GB in decimal
units, or about \(477.8\) GiB, of raw candidate bytes before filtering,
sorting, and merging.  A computation of \(a(15)\) remains feasible as a
larger segmented campaign, but it was outside the scope of the present run.

\section{Implementation and segmented search}
\label{sec:impl}

The production implementation is the C17 source file \path{src/a096242.c},
built by the accompanying \path{src/Makefile}.  It supports serial and
POSIX-threaded execution, in-memory and segmented modes, validation against
an OEIS-style b-file, and optional saving of each level as a sorted binary
file of little-endian \texttt{uint64\_t} values.  The program rejects
\texttt{--max-n} values above \(20\), since \(9^{21}-1\) exceeds the range
of the \texttt{uint64\_t} type.

For smaller levels the program generates the whole candidate level in
memory, sorts it, deduplicates it, and then tests primality.  For larger
levels it uses segmented generation: parent blocks are processed into sorted
temporary runs, each run is locally deduplicated, and a global k-way merge
removes duplicates spanning different runs while streaming candidates into
the primality stage.  The merge maintains a min-heap containing the current
value from each run; when successive heap outputs have the same value, only
the first occurrence is forwarded to the primality stage.  The segmented mode
changes storage strategy, not the mathematical set being enumerated.

In schematic form, each level is computed as follows:
\begin{enumerate}
\item generate insertions from \(D_{n-1}\);
\item reject leading-zero insertions and final digits \(0\), \(3\), and \(6\);
\item sort and globally deduplicate the surviving candidates;
\item apply small-prime trial division and deterministic Miller--Rabin;
\item write the accepted primes as the sorted level \(D_n\).
\end{enumerate}
In segmented mode, the sort-and-deduplicate step is implemented by sorted
temporary runs followed by a k-way merge; duplicates spanning runs are
removed during that merge.

The project hardware target recorded in the local dossier was a MinisForum
UM790 Pro with an AMD Ryzen 9 7940HS processor, 64 GB DDR5 memory, and 16
hardware threads.  This hardware information is included to document the
local computation environment, not as a hardware-independent performance
claim.

Exact operating-system, kernel, C library, and compiler-version strings for
the production run are not present in the local run-evaluation reports
consulted for this manuscript; they are therefore not stated here as
production metadata.

The release build in the current Makefile uses C17, POSIX threads, strict
warning flags, and optimized native compilation:
\begin{verbatim}
-std=c17 -Wall -Wextra -Wpedantic -Wshadow -Wconversion
-Wdouble-promotion -Wformat=2 -pthread
-O3 -march=native -flto -DNDEBUG
\end{verbatim}
The code-review artifacts also record clean GCC and Clang warning-matrix
builds.  Full wall-clock timings and peak RSS values for the original
production run that produced the saved \(D_{11}\) through \(D_{14}\) levels
are not included in the local run-evaluation reports; they are therefore not
reported here.  The local artifacts consulted here do not record a separate
machine identity for the \(D_{11}\) validation reruns.  The \(D_{11}\)
timings reported below are validation-rerun measurements, not recovered
original-production timings or hardware-independent benchmarks.

\section{Validation and reproducibility}
\label{sec:validation}

Validation was organized around the production source, the extended b-file
\path{data/b096242.txt}, the run-evaluation reports for \(n=13\) and
\(n=14\), the manifest \path{levels_manifest.tsv}, and the separate
saved-level certifier \path{cert/a096242_level_certifier.c}.

Throughout this section, paths of the form \path{levels/DNN.bin} refer to
saved binary level files produced during the local computational campaign.
Those multi-GB generated artifacts are not included in the minimal public
GitHub package.  Their counts, byte sizes, and SHA-256 hashes are documented
in \path{levels_manifest.tsv} and in the validation reports.

The known OEIS values through \(n=10\) were replayed by serial, parallel
in-memory, and forced segmented runs, with counts matching the local OEIS
snapshot as recorded in the local validation summary.  The code-review
artifacts record release, debug, and sanitizer builds; regression tests;
deterministic comparison of saved levels for different thread counts through
\(D_{10}\); ASan/UBSan, TSan, and Valgrind runs; and static-analysis checks.

After the manuscript review process identified the absence of a high-level
cross-mode check, the first new level \(D_{11}\) was regenerated by the
production program in two deliberately different modes: single-thread
in-memory mode and 16-thread segmented mode with parent chunks of size
10000.  Both regenerated files were byte-identical to the saved
\path{levels/D11.bin} file.  The shared SHA-256 hash was
\[
\texttt{38ffc99bf4836f55...af0dd3a59}.
\]
The hash is abbreviated in the display; full 64-character SHA-256 values are
recorded in \path{levels_manifest.tsv}.
The single-thread in-memory rerun reported \(a(11)=5056361\), wall time
0:19.48, and peak RSS 1718048 kB.  The 16-thread segmented rerun reported
\(a(11)=5056361\), wall time 0:10.56, and peak RSS 102848 kB.

The first new level was also re-enumerated by a separate C validation
program, \path{cert/a096242_d11_gmp_reenumerate.c}.  This program reads
\path{levels/D10.bin}, enumerates all raw \(D_{11}\) insertion attempts,
applies the leading-zero and base-9 last-digit filters, globally
deduplicates the surviving candidates, applies the same small-prime
trial-division pre-filter through \(97\), and then calls GMP's
\texttt{mpz\_probab\_prime\_p} with repetition parameter \(25\) rather than
using the production Miller--Rabin routine.  The local GMP package version
was 6.3.0.  The run reported 112,876,929 raw insertion attempts, 87,478,862
globally unique candidates after the leading-zero and base-9 last-digit
filters and global deduplication, 14,288,558 GMP primality calls, and
5,056,361 accepted values.  The resulting \path{D11_gmp.bin} file was byte-identical
to \path{levels/D11.bin}, with the same SHA-256 hash above.  This is an
independent probabilistic primality cross-check, not a deterministic
primality certificate.  The GMP re-enumeration ran in 0:18.07 wall time with
peak RSS 1,625,204 kB.

These \(D_{11}\) checks strengthen the evidence for the first new term and
for the production pipeline's memory/segmented equivalence at a new level.
They are still not an exhaustive independent re-enumeration of
\(D_{12}\), \(D_{13}\), or \(D_{14}\).  For those levels, the present
validation consists of the production enumeration, manifest hashes, structural
file checks, and the saved-level validity certifier described below.

The saved binary files \path{levels/D01.bin} through \path{levels/D14.bin},
representing \(D_1\) through \(D_{14}\), were checked for file-size
consistency, strict increasing order, absence of duplicates, and membership
in the correct base-9 digit-length interval.  The manifest records the count,
byte size, and SHA-256 hash of each level.  Table~\ref{tab:manifest} lists
the four new levels.

\begin{table}[ht]
\centering
\small
\begin{tabular}{r r r l}
\toprule
Level & Count & Bytes & SHA-256 prefix\\
\midrule
11 & 5,056,361 & 40,450,888 & \texttt{38ffc99bf4836f55...af0dd3a59}\\
12 & 22,726,556 & 181,812,448 & \texttt{42c916e29431c05c...207c40a7c}\\
13 & 103,351,642 & 826,813,136 & \texttt{69ff460d9198fbe4...e95126595}\\
14 & 475,048,816 & 3,800,390,528 & \texttt{13d23872161dcb27...8ba91aa4}\\
\bottomrule
\end{tabular}
\caption{Manifest entries for the four new levels.  The full SHA-256 values
are stored in \texttt{levels\_manifest.tsv}.}
\label{tab:manifest}
\end{table}

A separate saved-level certifier was run through \(n=14\).  It loads the
saved levels, verifies \(D_1=\{2,3,5,7\}\), checks sorted order and
digit-length range, and for every saved value in \(D_2,\ldots,D_{14}\)
checks primality and the existence of at least one valid base-9 digit
deletion into the previous saved level.  The report records
\[
  \texttt{CERTIFIER RESULT: PASS through n=14}.
\]
This certifier supports saved-level validity.  It does not perform a second
exhaustive enumeration proving that no candidate is missing.  Since the
definition of a deletable prime requires the existence of at least one valid
deletion into the previous level, this parent-existence check is sufficient
for membership of a saved value, conditional on primality and on the previous
saved level.

The certifier is independent of the production program as a separate C
source file and saved-level traversal.  Its primality predicate uses the same
deterministic 64-bit Miller--Rabin witness set as the production code, so it
should not be read as an independent primality-method certificate such as an
ECPP proof or a GMP-based oracle.  Its independent contribution is the
separate pass over the saved files: range and order checks, duplicate
rejection, parent-deletion lookup, and per-value validation against the
previous saved level.

\section{Results}
\label{sec:results}

The proposed extended b-file is shown in Table~\ref{tab:terms}.  The first
ten values are the values in the local OEIS snapshot and were replayed by
the implementation.  The last four values are the new computed terms.

\begin{table}[H]
\centering
\begin{tabular}{rr}
\toprule
\(n\) & \(a(n)\)\\
\midrule
1 & 4\\
2 & 14\\
3 & 58\\
4 & 221\\
5 & 911\\
6 & 3,638\\
7 & 14,687\\
8 & 61,435\\
9 & 262,189\\
10 & 1,140,171\\
11 & 5,056,361\\
12 & 22,726,556\\
13 & 103,351,642\\
14 & 475,048,816\\
\bottomrule
\end{tabular}
\caption{Values of A096242 through \(n=14\).}
\label{tab:terms}
\end{table}

\begin{proof}[Proof of Computational Result~\ref{thm:main}]
The mathematical completeness argument is inductive.  The base level is
\(D_1=\{2,3,5,7\}\).  If the saved level \(D_{n-1}\) is complete, then
Lemma~\ref{lem:bottom-up} says that \(D_n\) is obtained by considering the
distinct valid insertions into \(D_{n-1}\) and retaining exactly the primes.
Proposition~\ref{prop:last-digit} gives a necessary last-digit rejection for
multi-digit candidates in base 9, and it removes no prime.  The production
implementation follows this induction level by level, with global sort/merge
deduplication before primality testing.

The first ten counts reproduce the local OEIS b-file.  The saved levels
\(D_{11}\) through \(D_{14}\) have the counts shown in
Table~\ref{tab:manifest}, and the full count/byte/hash records are stored in
\path{levels_manifest.tsv}.  The separate saved-level certifier verifies
every saved value through \(D_{14}\) as prime with a valid deletion into the
previous level.  In addition, \(D_{11}\) was reproduced byte-for-byte by
single-thread memory and 16-thread segmented production reruns and by a
separate GMP-based re-enumerator.  These checks strengthen the first new
level and the pipeline reproducibility evidence, while completeness of
\(D_{12}\) through \(D_{14}\) remains tied to the documented production
enumeration.  Subject to the correctness of that enumeration and validation
pipeline, the cardinalities of these four levels are therefore the four
stated values.
\end{proof}

\begin{proof}[Proof of Corollary~\ref{cor:bfile}]
The local OEIS source snapshot lists values through \(n=10\).
Computational Result~\ref{thm:main} supplies the next four values computed
by the present project.  Combining these two sources gives exactly the
fourteen rows shown in Table~\ref{tab:terms}.
\end{proof}

\section{Conclusion}
\label{sec:conclusion}

The bottom-up formulation turns the recursive definition of base-9 deletable
primes into an exact level-by-level enumeration.  Applied to A096242, it
extends the known table from \(n=10\) to \(n=14\), adding
\[
  5056361,\quad 22726556,\quad 103351642,\quad 475048816.
\]
The computation is best read as a finite certified table extension, not as a
growth law for deletable primes.  Its distinctive features are the exact
parent-frontier generation, the base-9 last-digit filter, the need for
global deduplication in a level DAG, and the manifest-based reproducibility
of large generated levels.

\section*{Data availability}

The code, extended b-file, validation reports, saved-level certifier, and
level manifest are available in the public GitHub repository and its Zenodo
archive \cite{cortiA096242Package2026}.  The project repository is
\url{https://github.com/carcorti/A096242}, and the archived release is
available at \url{https://doi.org/10.5281/zenodo.20581301}.  The principal
artifacts are the C17 source and Makefile, \path{data/b096242.txt},
\path{levels_manifest.tsv}, the run-evaluation reports, the saved-level
certifier source, and the GMP-based \(D_{11}\) re-enumerator source
\path{cert/a096242_d11_gmp_reenumerate.c}.  The binary files
\path{levels/D01.bin} through \path{levels/D14.bin} are generated local
artifacts, not intended contents of the minimal GitHub repository.  They are
regenerable from the source code and checkable against the manifest counts,
byte sizes, and SHA-256 hashes.  A local Markdown mirror of Sinclair's
witness table is included in the public package as
\path{references/sinclair_miller_rabin_sprp_bases.md}.

\section*{AI use disclosure}

This work used large language models as research assistants during the
workflow for the A096242 project.  Gemini was used in the initial triage of
the sequence and project feasibility.  An OpenAI Codex system, identified by
the working environment as GPT-5, was used during code development, code
review, validation planning, manuscript drafting, and manuscript revision.
Additional LLM systems identified as Claude, DeepSeek, Grok, Kimi, Qwen, and
Z.ai were used as non-author reviewers during code and manuscript review
rounds.

The AI-assisted work included mathematical analysis of the base-9 deletable
prime definition, formulation of the bottom-up insertion method, recognition
of the base-9 last-digit filter, C17 implementation support, review triage,
validation planning, interpretation of run reports, and drafting of this
manuscript.

Mathematical analysis.  The bottom-up completeness lemma, the last-digit
filter, and the DAG/deduplication interpretation were developed and written
with AI assistance.  The initial sequence triage used Gemini; the formulation
and manuscript presentation of the bottom-up lemma, DAG interpretation, and
validation boundaries were developed mainly in Codex-assisted sessions.  The
identities and set-theoretic claims used in the paper are proved explicitly
in the text and were checked against the final implementation and saved-level
certifier before inclusion.
The human author independently verified the logical correctness of each proof
and the validity of each algorithmic step before accepting the final text.

Algorithm and code.  The production C17 search program, Makefile,
saved-level certifier, and GMP-based \(D_{11}\) re-enumerator were produced
through an AI-assisted development and review workflow.  Their correctness
does not rest on that provenance: it rests on known-term replay,
memory/segmented-mode validation, sanitizer and tooling checks,
deterministic saved-level comparisons, run-evaluation reports, manifest
hashes, the separate saved-level validity check, and the independent
GMP-based \(D_{11}\) re-enumeration described in this paper.

Manuscript.  The manuscript was drafted with AI assistance and revised by the
author.  AI systems supplied drafts, edits, and review suggestions; the final
verifying revision and acceptance of the text were performed by the human
author.  Every quantitative claim included here was checked against the local
computational artifacts: the OEIS source snapshot, the extended b-file, the
level manifest, the run-evaluation reports, the validation summary, and the
certifier report.

The author takes full scientific responsibility for the mathematical
statements, the algorithm, the implementation, the numerical results, the
validation boundaries, and the conclusions of this paper.

\begin{thebibliography}{9}
\sloppy
\raggedright

\bibitem{OEIS-A096242}
OEIS Foundation Inc.,
\emph{Sequence A096242: Number of \(n\)-digit base-9 deletable primes},
The On-Line Encyclopedia of Integer Sequences,
\url{https://oeis.org/A096242}.
Sequence created by Michael Kleber, Feb. 28, 2003; five more terms from
Ryan Propper, Jul. 19, 2005; Mathematica program by Robert Price,
Nov. 13, 2018.  Local project snapshot used as source material; live entry
checked by the author on June 6, 2026.

\bibitem{WeissteinDeletable}
Eric W. Weisstein,
\emph{Deletable Prime},
MathWorld--A Wolfram Web Resource.
\url{https://mathworld.wolfram.com/DeletablePrime.html}.
Accessed June 6, 2026.

\bibitem{Miller1976}
Gary L. Miller,
\emph{Riemann's hypothesis and tests for primality},
Journal of Computer and System Sciences 13 (1976), 300--317.
doi:10.1016/S0022-0000(76)80043-8.

\bibitem{Rabin1980}
Michael O. Rabin,
\emph{Probabilistic algorithm for testing primality},
Journal of Number Theory 12 (1980), 128--138.
doi:10.1016/0022-314X(80)90084-0.

\bibitem{SinclairSPRP}
Jim Sinclair,
\emph{Strong pseudoprime bases},
online table of deterministic Miller--Rabin bases.
\url{https://miller-rabin.appspot.com/}.
Accessed June 6, 2026.  A local Markdown mirror of the cited table is
included in the public package as
\path{references/sinclair_miller_rabin_sprp_bases.md}.

\bibitem{Jaeschke1993}
Gerhard Jaeschke,
\emph{On strong pseudoprimes to several bases},
Mathematics of Computation 61 (1993), 915--926.
doi:10.1090/S0025-5718-1993-1192971-8.

\bibitem{SorensonWebster2017}
Jonathan P. Sorenson and Jonathan Webster,
\emph{Strong pseudoprimes to twelve prime bases},
Mathematics of Computation 86 (2017), no. 304, 985--1003.
doi:10.1090/mcom/3134.

\bibitem{ForisekJancina2015}
Michal Fori\v{s}ek and Jakub Jan\v{c}ina,
\emph{Fast primality testing for integers that fit into a machine word},
SOFSEM 2015 Student Research Forum Papers and Posters, CEUR Workshop
Proceedings 1326 (2015), 20--30.
\url{https://ceur-ws.org/Vol-1326/020-Forisek.pdf}.

\bibitem{cortiA096242Package2026}
C. Corti,
\emph{A096242: Computational package for OEIS sequence A096242},
GitHub repository and Zenodo archive, 2026.
\url{https://github.com/carcorti/A096242}.
\url{https://10.5281/zenodo.20581301}.

\end{thebibliography}

\end{document}
